\section{Proposed Method}
\label{sec:method}

\subsection{Overview}
\label{sec:overview}

ESC-Net~\cite{lee2024escnet} is a one-stage network that couples a frozen
CLIP~\cite{radford2021learning} encoder with the pre-trained decoder blocks of the
Segment Anything Model (SAM)~\cite{kirillov2023segment} through a vision--language
correlation map. \methodname{} adapts this paradigm from multi-class semantic
segmentation to the class-agnostic, binary problem of salient object detection
(SOD). \cref{fig:framework} shows the architecture. Given an input image $I$, a
CLIP ViT-L/14 encoder produces multi-scale visual features $F_v$, while the CLIP
text encoder embeds a \emph{bank} of $K$ salient-object descriptions into language
features $F_l$. From these we form a multi-prompt vision--language correlation
$C \in \mathbb{R}^{2 \times H \times W}$, which is progressively refined by $N$
stacked \emph{ESC blocks}, each combining a Pseudo Prompt Generator (PPG), a SAM
block, and a Vision--Language Fusion (VLF) block. \rev{Before decoding, the initial
prompt-derived correlation is added back to the refined correlation through a
residual connection (\cref{sec:residual}). A decoder then produces the saliency map
$M$. We study two variants: a simple configuration with a U-Net decoder trained with
binary cross-entropy, which produces all main results (\cref{tab:sota,tab:sip_rgbd}),
and a correlation-aware variant with an ASPP/FPN decoder and a boundary-aware
compound objective, used in the ablations (\cref{sec:decoder}).}
%% OLD: "Crucially, the initial prompt-derived correlation is preserved via a residual
%% connection ... A simple U-Net decoder then produces the saliency map M, and the
%% whole network is trained end-to-end with a boundary-aware compound objective."

\begin{figure*}[t]
  \centering
  \includegraphics[width=\linewidth]{SOD_framework.png}
  \caption{Overview of \methodname. A CLIP ViT-L/14 encoder provides multi-scale
    visual features, and the CLIP text encoder embeds a bank of $K$ salient-object
    prompts. \rev{The $K$ prompt correlations are aggregated (element-wise maximum in
    the simple configuration; learned per-pixel attention in the correlation-aware
    variant), refined by $N$ ESC blocks, and the initial correlation is added back
    through a residual before decoding.} \emph{Stars} mark the components we
    redesign relative to ESC-Net.}
  \label{fig:framework}
\end{figure*}


\subsection{Multi-Scale Vision--Language Features}
\label{sec:features}

We attach forward hooks to the transformer blocks $\{5, 11, 17, 23\}$ of a CLIP
ViT-L/14 encoder operating at $336 \times 336$, and reshape each block's patch
tokens into a spatial map of size $H \times W = 24 \times 24$. We denote the
deepest map by $F_v^{\textit{high}}$ (block $23$) and a mid-level map by
$F_v^{\textit{mid}}$ (block $11$); the remaining maps are kept as skip features for
the decoder.

To measure similarity against the language features, the visual maps are
projected from the ViT hidden dimension ($1024$) to the shared CLIP embedding
space ($768$) by a learnable matrix $W_v$ initialized from CLIP's own visual
projection and fine-tuned with the new modules. For a spatial location $i$,
\begin{equation}
  \hat{F}_v^{l}(i) = \frac{W_v^{\top} F_v^{l}(i)}{\lVert W_v^{\top} F_v^{l}(i) \rVert},
  \qquad l \in \{\textit{high}, \textit{mid}\}.
\end{equation}
Each text description is encoded by the CLIP text encoder and $\ell_2$-normalized
to obtain a language embedding in the same space. \rev{The text encoder is kept
frozen. The implementation used for the submitted version of this paper bypassed
the text transformer (it projected only the end-of-text token embedding), so every
prompt embedding depended only on the prompt's token length. We fixed this, and all
results in this version were retrained with the full CLIP text encoder.}


\subsection{Multi-Prompt Vision--Language Correlation}
\label{sec:corr}

ESC-Net forms its correlation map from the cosine similarity between the visual
features and a \emph{single} text embedding per class. For SOD this is
sub-optimal: the appearance of ``salient objects'' is extremely diverse, and any
one phrasing aligns with only part of that distribution, so a single prompt
frequently fails to activate small or atypical targets.

We instead maintain a \emph{prompt bank}
$\mathcal{T} = \{T^1, \dots, T^K\}$ of $K$ diverse salient-object descriptions
(\eg, ``a photo of a salient object'', ``the most visually distinctive object in
the scene'', ``a high contrast object against its background''). The bank is
encoded once into fixed embeddings $\{t_k\}_{k=1}^{K}$; no per-image caption or
external captioning model is used, so at test time the model consumes only the RGB
image with all text prompts fixed a priori. For each scale $l$ and each prompt we
compute a cosine-similarity map,
\begin{equation}
  C_k^{l}(i) = \big\langle \hat{F}_v^{l}(i),\; t_k \big\rangle ,
  \qquad k \in \{1, \dots, K\}.
\end{equation}

\rev{\subsubsection{Max-pooled aggregation.}
In the simple configuration the $K$ maps are aggregated by an element-wise maximum,
$\tilde{C}^{l}(i) = \max_{k} C_k^{l}(i)$, so a region is kept if any prompt
activates it.}

\subsubsection{Learned prompt-attention aggregation.}
\rev{As an alternative to the maximum,} we learn per-pixel attention weights that
determine how much each prompt contributes at each spatial location. A lightweight
convolutional network
$\mathcal{A}$ takes the $K$-channel correlation stack as input and produces
attention logits:
\begin{equation}
  \alpha_k^{l}(i) = \frac{\exp\bigl(\mathcal{A}(C^{l})_k(i)\bigr)}
                         {\sum_{j=1}^{K}\exp\bigl(\mathcal{A}(C^{l})_j(i)\bigr)},
  \qquad
  \tilde{C}^{l}(i) = \sum_{k=1}^{K} \alpha_k^{l}(i)\, C_k^{l}(i).
  \label{eq:attn_agg}
\end{equation}
The attention in \cref{eq:attn_agg} is strictly more expressive than max-pooling:
it can replicate max by concentrating all weight on the strongest prompt, but it
can also learn soft combinations that suppress noisy false positives at specific
locations. \rev{The attention network is trained only if the loss depends on the
aggregated correlation. This holds for the correlation-aware variant, where
attention and max-pooling are within about $0.003$ $S_m$ of each other
(\cref{tab:prompt}); in the simple
configuration the decoder does not read the correlation, so the attention weights
receive no gradient (\cref{sec:residual}).}
%% OLD: "Our ablation confirms that attention aggregation consistently outperforms
%% max-pooling, particularly on challenging datasets with cluttered backgrounds."

We then apply per-sample min--max normalization,
\begin{equation}
  C^{l} = \frac{\tilde{C}^{l} - \min \tilde{C}^{l}}
               {\max \tilde{C}^{l} - \min \tilde{C}^{l} + \epsilon},
\end{equation}
and stack the two scales into the correlation tensor
$C = [\,C^{\textit{high}};\, C^{\textit{mid}}\,] \in \mathbb{R}^{2 \times H \times W}$.


\subsection{Prompt-Residual Connection}
\label{sec:residual}

\begin{revpar}
Each VLF block (\cref{sec:sam}) ends with a text cross-attention in which the
visual tokens are queries and the text embeddings are keys and values. In the
multi-class setting of ESC-Net there is one text token per class, but in binary SOD
the fusion block receives a single caption embedding, so the attention has one key.
The softmax over a single key is identically one, so every query returns the same
value vector and the VLF output is independent of the correlation and the visual
features (changing either changes the output by exactly zero). The refined
correlation $C_N$ thus carries no gradient path back to the prompt aggregation or
the visual projection.

A residual connection restores such a path. Let $C_0$ denote the initial
multi-prompt correlation and $C_N$ the correlation after $N$ ESC blocks. The decoder
receives
\begin{equation}
  C_{\textit{final}} = C_N + C_0 .
  \label{eq:residual}
\end{equation}
There is no mixing coefficient. With the correlation-aware decoder, which reads
$C_{\textit{final}}$, the residual is the only route by which the visual projection
and the attention aggregation receive gradient (\cref{sec:residual_exp}). In a
matched comparison over three seeds, however, removing the residual does not reduce
accuracy (\cref{tab:prompt}). The simple decoder does not read the
correlation, so in the simple configuration the residual has no effect and the
prompt bank acts only through the discrete pseudo prompts passed to SAM. We report
the residual as an analysed design element, not as a source of accuracy gains.
\end{revpar}
%% OLD (submitted): "we observe that after N refinement iterations the VLF output
%% becomes entirely determined by its own parameters and the SAM-refined features,
%% erasing the spatial prior ... We address this with a simple but critical residual
%% connection ... C_final = C_N + \gamma C_0, where \gamma is a mixing coefficient ...
%% Our experiments show that without this residual, the multi-prompt bank provides
%% negligible improvement regardless of the aggregation method; with it, learned
%% prompt attention yields consistent gains across all benchmarks. We analyse the
%% choice of \gamma ... a light residual (\gamma=0.3) is both effective and stable."
%% (The code has no \gamma; the "\gamma" sweep in the submitted Table 6 varied the
%% point-classification loss weight.)


\subsection{Pseudo Prompt Generator with Adaptive Fallback}
\label{sec:ppg}

\begin{figure}[!t]
  \centering
  \includegraphics[width=0.9\linewidth]{PPG_figure.png}
  \caption{The proposed Pseudo Prompt Generator (PPG). From the correlation map
    $C$, a spatial softmax yields a probability map $P$, which is thresholded to a
    binary map $B$ and clustered into $N_o$ regions by $k$-means; the Hadamard
    product $P \odot$ clusters localizes the per-region maxima, giving $N_o$
    pseudo points $X_{\textit{pseudo}}$ and pseudo masks $M_{\textit{pseudo}}$ that
    SAM's prompt encoder maps to sparse and dense prompts. \emph{Star} marks
    our addition: the \emph{adaptive fallback}, which selects the top-$\rho$
    fraction of $P$ whenever the binary support $|B|$ is below $\tau$.}
  \label{fig:ppg}
\end{figure}

The PPG converts a correlation map into SAM-compatible point and mask prompts
(\cref{fig:ppg}). Following ESC-Net, we first apply a spatial softmax to the
(high-level) correlation map to obtain an object-probability map $P$, binarize it with a
threshold $\alpha$, and cluster the active pixels into $N_o$ regions by $k$-means
on pixel coordinates~\cite{hartigan1979algorithm}. For each region we emit one
\emph{pseudo point}, placed at the pixel of maximum probability and assigned a
positive label, and one \emph{pseudo mask} given by the region support. The $N_o$
points and masks are embedded by SAM's prompt encoder into sparse and dense
prompt features.

A direct port of this procedure is unstable for binary SOD. After a softmax over
$H\!\times\!W$ locations the mean probability is $\approx 1/(HW)$, so a fixed
threshold demands tens of times the average activation; for diffuse correlation
maps, which are common early in training and for small or low-contrast salient
objects, the binary mask is empty and no prompt is produced. We therefore add an
\emph{adaptive fallback}: whenever the active support falls below a minimum count
$\tau$, we instead select the top-$\rho$ fraction of $P$,
\begin{equation}
  B =
  \begin{cases}
    \mathbb{1}\!\left[P > \alpha\right], & \text{if } \big\lvert \mathbb{1}[P > \alpha] \big\rvert \geq \tau, \\[4pt]
    \mathrm{TopK}_{\rho}(P),            & \text{otherwise,}
  \end{cases}
\end{equation}
where $\tau = \max(4 N_o, \lfloor \rho H W \rfloor)$. Because the softmax always
has a positive maximum, $\mathrm{TopK}_{\rho}(P)$ is never empty, so the PPG is
\emph{guaranteed} to deliver substantive spatial guidance to the SAM block in
every iteration.


\subsection{SAM Block and Vision--Language Fusion}
\label{sec:sam}

The SAM block aggregates spatial context into the visual features using the
pseudo prompts, reusing the pre-trained transformer decoder blocks of SAM so that
their learned image--prompt mapping is preserved. \rev{The simple configuration uses
image-to-token cross-attention only; restoring the token-to-image step of SAM's
decoder layout (bidirectional attention, BiSAM) is evaluated in the ablation.}
%% OLD: "We follow the standard SAM decoder layout with bidirectional cross-attention
%% between prompt tokens and image features."

The VLF block refines the correlation map using the SAM-refined features and
the language embeddings through Swin-style window attention~\cite{liu2021swin}
followed by text-guided cross-attention, following CAT-Seg~\cite{cho2024cat}
and ESC-Net~\cite{lee2024escnet}. After $N$ ESC blocks the prompt-residual
(\cref{eq:residual}) is applied before decoding. \rev{As discussed in
\cref{sec:residual}, with a single caption key the text cross-attention makes the
VLF output input-independent. Our released code also provides an optional variant
in which the VLF attends over the $K$ bank prompts plus the caption ($K{+}1$ keys)
and adds its output to the incoming correlation; all results reported here use the
original block.}
%% OLD: "so that the segmentation loss can back-propagate through the initial prompt
%% correlation C_0 as well as the refined C_N."


\subsection{Decoder and Training Objective}
\label{sec:decoder}

\rev{The simple configuration uses a U-Net-style decoder~\cite{ronneberger2015u}
with bilinear upsampling and skip connections from the CLIP feature maps; it does
not take the correlation as input. The correlation-aware variant (``Dec'' in the
ablation) uses an ASPP/FPN decoder~\cite{chen2017rethinking,lin2017feature} that
injects $C_{\textit{final}}$ as an explicit saliency prior before the final
prediction layer.}
%% OLD: "Before the final prediction layer we inject the residual-enhanced correlation
%% C_final as an explicit saliency prior. Our ablation shows that a simple decoder
%% with the prompt-residual outperforms a complex ASPP/FPN decoder without it."

\subsubsection{Training objective.}
\rev{The simple configuration is trained with binary cross-entropy. The
correlation-aware variant} is trained with a boundary-aware compound loss
combining Tversky loss~\cite{salehi2017tversky} and Sobel-edge-weighted focal
BCE~\cite{lin2017focal}, plus a small probability-space MSE term:
\begin{equation}
  \mathcal{L} =
    w_{\textit{tv}}\, \mathcal{L}_{\textit{tv}} +
    w_{f}\, \mathcal{L}_{\textit{focal}} +
    \lambda\, \mathcal{L}_{\textit{mse}},
\end{equation}
with $w_{\textit{tv}}{=}0.6$, $w_{f}{=}0.4$, and $\lambda{=}0.1$.
Deep supervision is applied at two intermediate resolutions with weight $0.3$
\rev{(correlation-aware decoder only). The ablation additionally evaluates an
auxiliary point-classification loss with weight $\gamma_{\textit{pt}}$, which
compares the predicted saliency at each pseudo-point location with the ground
truth.}
